\section{Results}
Table~\ref{tab:main-results} summarizes WER and latency across the three evaluation conditions. CBCI
outperforms both the causal baseline and the uniform look-ahead system at matched latency in every
condition, and recovers the majority of the gap to the offline upper bound.

\begin{table}[t]
\centering
\caption{Word error rate (\%) and algorithmic latency across evaluation conditions on \textsc{Meridian-Conv}. Lower WER is better.}
\label{tab:main-results}
\footnotesize
\begin{tabular}{@{}lccc@{}}
\toprule
\textbf{System} & \textbf{Clean} & \textbf{Reverb} & \textbf{Babble 5dB} \\
\midrule
Causal baseline        & 8.9  & 11.4 & 15.2 \\
Uniform look-ahead      & 8.1  & 10.6 & 14.0 \\
\textbf{CBCI (ours)}    & \textbf{7.4} & \textbf{9.6}  & \textbf{12.7} \\
\midrule
Offline (upper bound)  & 6.2  & 8.0  & 10.5 \\
\bottomrule
\end{tabular}
\end{table}

Figure~\ref{fig:latency-tradeoff} plots WER against algorithmic latency as the look-ahead buffer $B$ is
varied at inference time for a single CBCI checkpoint, compared against uniform look-ahead models that
each require separate training runs. The buffer-annealed CBCI checkpoint traces a strictly better
frontier, particularly in the 100--250\,ms latency band relevant to live captioning.

\begin{figure}[t]
\centering
\begin{tikzpicture}
\begin{axis}[
  width=0.95\linewidth,
  height=4.6cm,
  xlabel={Algorithmic latency (ms)},
  ylabel={WER (\%)},
  xlabel style={font=\footnotesize},
  ylabel style={font=\footnotesize},
  tick label style={font=\footnotesize},
  legend style={font=\footnotesize,at={(0.97,0.95)},anchor=north east,draw=isgray!40},
  grid=major,
  grid style={isgray!15},
  axis lines=left,
  ymin=6, ymax=12,
]
\addplot[color=isblue,mark=*,thick] coordinates {
  (0,8.9) (60,8.3) (120,7.8) (160,7.4) (240,7.0) (320,6.8)
};
\addlegendentry{CBCI (single ckpt.)}
\addplot[color=isamber,mark=square*,thick,dashed] coordinates {
  (0,8.9) (60,8.6) (120,8.3) (160,8.1) (240,7.7) (320,7.4)
};
\addlegendentry{Uniform look-ahead}
\end{axis}
\end{tikzpicture}
\caption{Latency-accuracy trade-off on the clean condition. Each uniform look-ahead point requires a
separately trained model; CBCI sweeps $B$ at inference time from one checkpoint.}
\label{fig:latency-tradeoff}
\end{figure}

\subsection{Ablations}
Removing buffer annealing (training with fixed $B{=}160$\,ms only) degrades WER by 0.4 absolute when
evaluated at $B{=}80$\,ms, confirming that annealing is responsible for graceful degradation at unseen
buffer sizes. Restricting $\mathcal{S}$ to a single layer recovers only 38\% of the full CBCI gain,
while extending $\mathcal{S}$ to all 18 layers matches the 6-layer configuration on WER but doubles peak
memory during training, indicating diminishing returns beyond the greedily selected set.
